Transformer-based models \citep{NIPS2017_3f5ee243} have become foundational across diverse domains, including computer vision \citep{dosovitskiy2021an,bao2022beit,he2022masked,peebles2023scalable}, natural language \citep{devlin-etal-2019-bert,Lan2020ALBERT, brown2020language}, robotics \citep{brohan2022rt,chen2021decision}, and scientific discovery \citep{jumper2021highly, kamienny2022end}. In recent years, there has been evidence that scaling up model size, dataset size, or computational budget during pre-training can yield unprecedented performance gains \citep{kaplan2020scalinglawsneurallanguage}, driving the proliferation of Transformer-based foundation models \citep{openai2024gpt4technicalreport,dubey2024llama,anil2023gemini,oquab2024dinov}. 

Despite these achievements, the architecture of Transformers—especially the role of individual layers—remains largely heuristic. For instance, empirical studies have observed high redundancy in the deeper layers \citep{gromov2025the,men-etal-2025-shortgpt}, uniformity of representations in the middle layers \citep{sun2025transformer}, and robustness to permuting certain intermediate layers \citep{lad2025remarkable} in LLMs. These findings suggest convergent functionality that one layer represents, yet our understanding of its role in processing information and representation learning remains limited. While interpretability efforts to unveil the function underlying network layers, especially Transformer blocks, exist---spanning mechanistic interpretability \citep{elhage2021mathematical,nanda2023progress,wang2023interpretability,huben2024sparse}, causal mediation analysis \citep{vig2020investigating,meng2022locating}, and visualization \citep{bricken2023monosemanticity,olsson2022context}---most of them focus on \emph{post-hoc} interpretation and phenomenological approaches. This motivates a pivotal question: 
% \emph{Can we find or design a function prior that induces a model that is interpretable by construction?}
\begin{tcolorbox}
[highlightquestion, boxrule=1pt]
    \centering
    \emph{Can we find or design a function prior that induces a model interpretable by construction?}
\end{tcolorbox}

One approach to achieving intrinsic interpretability is to embed an explicit optimization process into neural networks, known as model-based deep learning \citep{shlezinger2023model}. Prior works have designed networks that solve domain-specific problems such as constraint satisfaction \citep{wang2019satnet}, optimal control \citep{amos2017optnet, amos2018differentiable}, or physical simulation \citep{greydanus2019hamiltonian,karniadakis2021physics}. However, these models often rely on fixed task priors and lack generality.

Another more general avenue is energy-based learning (EBL) \citep{dawid2024introduction}, which frames prediction as minimizing a scalar energy function $E_\theta(x, y)$ over outputs $y$ conditioned on inputs $x$. Within this framework, Energy Transformer \citep{hoover2024energy} interprets Transformer layers as iterative optimization over the canonical continuous Hopfield energy \citep{ramsauer2021hopfield,krotov2021large} yet focuses on mechanistic analogies to associative memory without grounding its formulation in specific representational challenges.
In contrast, our goal is to find a design principle from top down that can not only reinterpret existing components but is also generalizable to novel architecture design \emph{constructively}.

In this work, we therefore take a fundamentally different approach by introducing a principle grounded in maximum likelihood estimation (MLE) for tokens on the hypersphere. Under mild assumptions, we interpret it under two complementary objectives for representation dynamics: semantic alignment (mode seeking) in high-dimensional space and distributional uniformity (mass covering) in a low-dimensional subspace. To translate these objectives into optimizable quantities over tokens, we define two complementary Hopfield-style energy functions that quantify these objectives and can be minimized through iterative optimization. This leads to the \emph{Hyper-Spherical Energy Transformer} (\textsc{Hyper-SET})—a recurrent-depth model in which core components such as symmetric attention, feedforward layers, $\operatorname{RMSNorm}$, and skip connection emerge naturally from the optimization dynamics. With only one set of shared parameters across iterations, \textsc{Hyper-SET} is compact, interpretable by design, and empirically competitive across diverse tasks, including reasoning, classification, and masked image modeling. Beyond a single instantiation, this principle can induce novel architectural designs by generalizing the energy functions, enabling variants such as linear attention and gated feedforward layer. To enhance scalability, we introduce depth-wise low-rank adaptation (LoRA), allowing flexible iteration-specific modulation with minimal parameter overhead. Our key contributions are summarized as follows:
\begin{enumerate}[leftmargin=*]
\item \textbf{Theoretical Formulation}: We conceptualize a general principle for information processing in layer dynamics based on maximum likelihood estimation on the hypersphere with two properties: uniformity and alignment, quantified via complementary Hopfield-style energy functions.
\item \textbf{Energy-Driven Architecture}: We derive a compact Transformer-based model through pure energy minimization, where core components—including symmetric attention, feedforward, $\operatorname{RMSNorm}$~\citep{zhang2019root}, and skip connection—emerge naturally.
\item \textbf{Competitive Performance}: We show competitive performance to vanilla Transformer across reasoning, classification, and masked modeling while demonstrating generality to design novel components (\eg, linear attention, gated feedforward) and scalability with flexible computation.
\end{enumerate}